% Esquema del loop de cálculo iterativo de la potencia límite por DNBR.
% Incluir en el documento principal con, por ejemplo:
%   \begin{figure}[h]
%       \centering
%       % Esquema del loop de cálculo iterativo de la potencia límite por DNBR.
% Incluir en el documento principal con, por ejemplo:
%   \begin{figure}[h]
%       \centering
%       % Esquema del loop de cálculo iterativo de la potencia límite por DNBR.
% Incluir en el documento principal con, por ejemplo:
%   \begin{figure}[h]
%       \centering
%       % Esquema del loop de cálculo iterativo de la potencia límite por DNBR.
% Incluir en el documento principal con, por ejemplo:
%   \begin{figure}[h]
%       \centering
%       \input{Figures/loop_potencia_dnbr.tex}
%       \caption{Esquema del loop de cálculo de la potencia límite por DNBR.}
%       \label{fig:loop_potencia_dnbr}
%   \end{figure}
\begin{tikzpicture}[
    font=\small,
    >={Latex[length=2mm]},
    node distance=10mm and 22mm,
    box/.style={
        rectangle, rounded corners, draw, thick,
        fill=blue!8, align=center,
        text width=60mm, minimum height=6mm,
        inner sep=1mm
    },
    startbox/.style={
        box, fill=gray!15
    },
    decision/.style={
        diamond, draw, thick, fill=orange!12, aspect=2.2,
        align=center, text width=30mm, inner sep=1mm
    },
    endbox/.style={
        rectangle, rounded corners, draw, thick,
        fill=green!12, align=center,
        text width=20mm, minimum height=8mm, inner sep=1mm
    },
    arrow/.style={-, draw, thick},
]

    % Nodos principales
    \node[startbox]                        (init)  {Inicialización de potencia con semilla};
    \node[box, below=of init]              (caudal){Cálculo de caudal másico};
    \node[box, below=of caudal]            (cobra) {Correr COBRA3-CP};
    \node[box, below=of cobra]             (extrae){Extraer potencia para DNBR límite};
    \node[decision, below=of extrae]       (conv)  {¿Convergencia en caudal y potencia?};
    \node[endbox, right=28mm of conv]      (fin)   {Fin};

    % Flujo principal
    \draw[->, thick] (init)   -- (caudal);
    \draw[->, thick] (caudal) -- (cobra);
    \draw[->, thick] (cobra)  -- (extrae);
    \draw[->, thick] (extrae) -- (conv);
    \draw[->, thick] (conv)   -- node[above] {Sí} (fin);

    % Loop de realimentación: si no converge, se vuelve al paso 2
    \coordinate (loopX) at ($(caudal.west) + (-15mm,0)$);
    \draw[->, thick] (conv.west) -- (loopX |- conv.west) node[pos=0.35, above] {No}
        -- (loopX) -- (caudal.west);

\end{tikzpicture}

%       \caption{Esquema del loop de cálculo de la potencia límite por DNBR.}
%       \label{fig:loop_potencia_dnbr}
%   \end{figure}
\begin{tikzpicture}[
    font=\small,
    >={Latex[length=2mm]},
    node distance=10mm and 22mm,
    box/.style={
        rectangle, rounded corners, draw, thick,
        fill=blue!8, align=center,
        text width=60mm, minimum height=6mm,
        inner sep=1mm
    },
    startbox/.style={
        box, fill=gray!15
    },
    decision/.style={
        diamond, draw, thick, fill=orange!12, aspect=2.2,
        align=center, text width=30mm, inner sep=1mm
    },
    endbox/.style={
        rectangle, rounded corners, draw, thick,
        fill=green!12, align=center,
        text width=20mm, minimum height=8mm, inner sep=1mm
    },
    arrow/.style={-, draw, thick},
]

    % Nodos principales
    \node[startbox]                        (init)  {Inicialización de potencia con semilla};
    \node[box, below=of init]              (caudal){Cálculo de caudal másico};
    \node[box, below=of caudal]            (cobra) {Correr COBRA3-CP};
    \node[box, below=of cobra]             (extrae){Extraer potencia para DNBR límite};
    \node[decision, below=of extrae]       (conv)  {¿Convergencia en caudal y potencia?};
    \node[endbox, right=28mm of conv]      (fin)   {Fin};

    % Flujo principal
    \draw[->, thick] (init)   -- (caudal);
    \draw[->, thick] (caudal) -- (cobra);
    \draw[->, thick] (cobra)  -- (extrae);
    \draw[->, thick] (extrae) -- (conv);
    \draw[->, thick] (conv)   -- node[above] {Sí} (fin);

    % Loop de realimentación: si no converge, se vuelve al paso 2
    \coordinate (loopX) at ($(caudal.west) + (-15mm,0)$);
    \draw[->, thick] (conv.west) -- (loopX |- conv.west) node[pos=0.35, above] {No}
        -- (loopX) -- (caudal.west);

\end{tikzpicture}

%       \caption{Esquema del loop de cálculo de la potencia límite por DNBR.}
%       \label{fig:loop_potencia_dnbr}
%   \end{figure}
\begin{tikzpicture}[
    font=\small,
    >={Latex[length=2mm]},
    node distance=10mm and 22mm,
    box/.style={
        rectangle, rounded corners, draw, thick,
        fill=blue!8, align=center,
        text width=60mm, minimum height=6mm,
        inner sep=1mm
    },
    startbox/.style={
        box, fill=gray!15
    },
    decision/.style={
        diamond, draw, thick, fill=orange!12, aspect=2.2,
        align=center, text width=30mm, inner sep=1mm
    },
    endbox/.style={
        rectangle, rounded corners, draw, thick,
        fill=green!12, align=center,
        text width=20mm, minimum height=8mm, inner sep=1mm
    },
    arrow/.style={-, draw, thick},
]

    % Nodos principales
    \node[startbox]                        (init)  {Inicialización de potencia con semilla};
    \node[box, below=of init]              (caudal){Cálculo de caudal másico};
    \node[box, below=of caudal]            (cobra) {Correr COBRA3-CP};
    \node[box, below=of cobra]             (extrae){Extraer potencia para DNBR límite};
    \node[decision, below=of extrae]       (conv)  {¿Convergencia en caudal y potencia?};
    \node[endbox, right=28mm of conv]      (fin)   {Fin};

    % Flujo principal
    \draw[->, thick] (init)   -- (caudal);
    \draw[->, thick] (caudal) -- (cobra);
    \draw[->, thick] (cobra)  -- (extrae);
    \draw[->, thick] (extrae) -- (conv);
    \draw[->, thick] (conv)   -- node[above] {Sí} (fin);

    % Loop de realimentación: si no converge, se vuelve al paso 2
    \coordinate (loopX) at ($(caudal.west) + (-15mm,0)$);
    \draw[->, thick] (conv.west) -- (loopX |- conv.west) node[pos=0.35, above] {No}
        -- (loopX) -- (caudal.west);

\end{tikzpicture}

%       \caption{Esquema del loop de cálculo de la potencia límite por DNBR.}
%       \label{fig:loop_potencia_dnbr}
%   \end{figure}
\begin{tikzpicture}[
    font=\small,
    >={Latex[length=2mm]},
    node distance=10mm and 22mm,
    box/.style={
        rectangle, rounded corners, draw, thick,
        fill=blue!8, align=center,
        text width=60mm, minimum height=6mm,
        inner sep=1mm
    },
    startbox/.style={
        box, fill=gray!15
    },
    decision/.style={
        diamond, draw, thick, fill=orange!12, aspect=2.2,
        align=center, text width=30mm, inner sep=1mm
    },
    endbox/.style={
        rectangle, rounded corners, draw, thick,
        fill=green!12, align=center,
        text width=20mm, minimum height=8mm, inner sep=1mm
    },
    arrow/.style={-, draw, thick},
]

    % Nodos principales
    \node[startbox]                        (init)  {Inicialización de potencia con semilla};
    \node[box, below=of init]              (caudal){Cálculo de caudal másico};
    \node[box, below=of caudal]            (cobra) {Correr COBRA3-CP};
    \node[box, below=of cobra]             (extrae){Extraer potencia para DNBR límite};
    \node[decision, below=of extrae]       (conv)  {¿Convergencia en caudal y potencia?};
    \node[endbox, right=28mm of conv]      (fin)   {Fin};

    % Flujo principal
    \draw[->, thick] (init)   -- (caudal);
    \draw[->, thick] (caudal) -- (cobra);
    \draw[->, thick] (cobra)  -- (extrae);
    \draw[->, thick] (extrae) -- (conv);
    \draw[->, thick] (conv)   -- node[above] {Sí} (fin);

    % Loop de realimentación: si no converge, se vuelve al paso 2
    \coordinate (loopX) at ($(caudal.west) + (-15mm,0)$);
    \draw[->, thick] (conv.west) -- (loopX |- conv.west) node[pos=0.35, above] {No}
        -- (loopX) -- (caudal.west);

\end{tikzpicture}
